\title{Direct Preference Optimization: Your Language Model is Secretly a Reward Model}

\begin{abstract}
Large-scale unsupervised language models (LMs) are effective at acquiring extensive world knowledge and some reasoning capabilities, yet exerting precise influence over their output remains challenging due to the absence of supervision in their training process. To introduce controllability, current approaches typically collect human-annotated rankings of generated responses and subsequently adapt the pretrained LM to follow these preferences, most often via reinforcement learning from human feedback (RLHF). Nonetheless, RLHF involves a multi-stage and frequently unstable approach: it requires first training a reward model that encodes human preferences, and then leveraging reinforcement learning to optimize the LM towards this reward, all while preventing excessive deviation from the original pretrained model. This work proposes a novel way of expressing the reward model in RLHF, making it possible to derive the associated optimal policy in closed form. This reformulation allows the conventional RLHF objective to be solved using a straightforward classification loss. Our proposed algorithm, termed \textit{Direct Preference Optimization} (DPO), is efficient, robust, and simple, dispensing with the necessity for generation sampling during fine-tuning and alleviating extensive hyperparameter tuning. Empirical results indicate that DPO can align LMs to human preferences as well as, or better than, prevailing alternatives. In particular, DPO achieves superior sentiment control compared to PPO-based RLHF, and equals or surpasses these methods in summarization and single-turn dialogue tasks, all while being considerably easier to deploy and train.
\end{abstract}

\section{Introduction}
Large-scale unsupervised language models (LMs) trained on extensive datasets have demonstrated unexpected abilities~\citep{chowdhery2022palm, brown2020language, touvron2023llama, bubeck2023sparks}. Nevertheless, these models learn from human-generated data that reflects a broad spectrum of objectives, values, and proficiencies. Not all of these are appropriate to replicate; for instance, an AI coding assistant should ideally \textit{recognize} frequent programming errors in order to correct them, but when writing code, it should prioritize the often less-common, high-caliber coding examples embedded in its training data. Likewise, while a language model ought to be \textit{informed} about prevalent misconceptions held by, say, half the population, we do not want the model to affirm such misunderstandings as factual in half of all related queries. The essential challenge is therefore to guide the model to express its \emph{preferred outputs and conduct} out of its vast reservoir of \textit{knowledge and competencies}, a necessity for crafting AI systems that are safe, effective, and controllable \citep{ouyang2022training}. Standard approaches tend to use reinforcement learning (RL) to adapt LMs towards aligning with human preferences; we demonstrate that the conventional RL-based training goal can actually be solved directly using a straightforward binary cross-entropy loss, offering a significant simplification to the process of preference alignment.

Broadly speaking, prevalent techniques enforce desired behaviors in a language model through carefully curated collections of human preference data, representing the behavioral traits considered beneficial and safe. This stage of preference alignment is typically implemented subsequent to the initial phase of large-scale unsupervised pre-training on text. While a direct approach involves supervised fine-tuning using annotated, high-quality human responses, the dominant and most effective paradigm is reinforcement learning from human (or AI) feedback (RLHF/RLAIF; \citep{christiano2017deep, bai2022constitutional}). RLHF methodologies involve training a reward model on human preference annotations and subsequently using RL to adapt a language model policy, enhancing the likelihood of responses that receive higher reward while preventing significant deviation from the base model. Although RLHF has achieved remarkable progress in conversational and coding proficiency, its training pipeline is far more intricate than that of supervised learning, requiring the training of several LMs and on-policy sampling within the training loop, thereby imposing notable computational demands.

This work presents a method for aligning language models with human preferences by optimizing them directly, bypassing both explicit reward modeling and reinforcement learning. We introduce \textit{Direct Preference Optimization (DPO)}, an approach that, while optimizing the same objective as prevalent RLHF algorithms (maximizing reward subject to a KL-divergence constraint), offers a more accessible implementation and more straightforward training procedure. At a high level, DPO adjusts the model by boosting the log probability of responses favored by humans over those that are not, leveraging a per-example, adaptive importance weight. This weighting is crucial in avoiding model collapse—a problem observed when employing a simple probability ratio objective. Like related techniques, DPO utilizes a theoretical model of preference, such as the Bradley-Terry model \cite{bradley1952rankanalysis}, which quantifies how closely a reward function aligns with actual preference data. In contrast to other methods that utilize the preference model to construct a loss for training a reward model, followed by policy optimization to maximize this learned reward, DPO applies a variable transformation to directly formulate the preference loss in terms of the policy itself. With a dataset of human-labeled preferences over generated responses, DPO enables straightforward policy optimization using a binary cross entropy loss, leading to a policy that is optimal with respect to an implicit reward shaped by the empirical preferences.

The principal contribution of our study is Direct Preference Optimization (DPO): a reinforcement learning-free, straightforward approach for training language models with preference data. Our empirical findings indicate that DPO performs on par with established methods—such as PPO-based RLHF—in tasks including sentiment adjustment, summarization, and conversational agents, and does so effectively even for models containing up to 6B parameters.

\section{Related Work}

As self-supervised language models have grown in scale, their capacity for zero-shot task completion \citep{radford2019language} and few-shot prompting \citep{gpt3,megatron,chowdhery2022palm} has similarly advanced. Nevertheless, these models often see notable enhancements in both downstream task performance and alignment with user goals after being fine-tuned on curated datasets comprising explicit instructions and corresponding human-authored completions \citep{mishra-etal-2022-cross,sanh2022multitask,chung2022scaling,thoppilan2022lamda}. This process, known as `instruction-tuning', enables large language models (LLMs) to extrapolate their understanding to novel instructions outside their training data, resulting in improved versatility and general effectiveness \citep{chung2022scaling}. Although instruction tuning has yielded promising results, collecting \textit{relative} human assessments of response quality is typically less resource-intensive than gathering expert-level demonstrations. As a consequence, more recent approaches have employed datasets that encode human preference information to further refine LLMs, leading to advancements in domains such as translation \citep{kreutzer-etal-2018-reliability}, summarization \citep{stiennon2022learning,ziegler2020finetuning}, creative writing \citep{ziegler2020finetuning}, and adherence to instructions \citep{ouyang2022training,ramamurthy2023is}. These strategies initially train a reward model—often through neural networks—on the basis of human preferences represented via frameworks like the Bradley-Terry model \citep{bradley1952rankanalysis}. Subsequently, the LLM is fine-tuned to optimize this reward signal, employing reinforcement learning algorithms such as REINFORCE \citep{williams1992reinforce}, proximal policy optimization (PPO; \cite{schulman2017proximal}), or related methodologies \citep{ramamurthy2023is}. Parallel research further utilizes LLMs, which have already been fine-tuned using human feedback, to generate synthetic datasets reflecting preferences on specific attributes—such as safety or harmlessness—leveraging limited direct supervision from humans, often provided only as textual rubrics for annotation \citep{bai2022constitutional}. Collectively, these developments unite research from two fronts: the application of reinforcement learning to train language models for various objectives~\citep{Ranzato2015SequenceLT,paulus2018a,wu2018learning}, and the broader investigation into methods for learning from human preferences \citep{christiano2017deep,kupcsik2018learning}. While utilizing relative human judgments holds considerable promise, the process of reinforcement learning-based fine-tuning of large language models remains a significant logistical hurdle; this paper introduces a theoretically principled alternative for optimizing relative preferences without resorting to reinforcement learning.

Beyond applications in language, the task of inferring policies from preferences has received attention in both bandit and reinforcement learning frameworks, leading to a range of methodologies. When contextual bandits are trained on preferences or rankings instead of direct reward signals, this is formalized as a contextual dueling bandit (CDB; \cite{yue2012karmed,dudik2015contextual}). In settings where reward values are not available, the analysis of CDBs shifts from identifying an optimal policy to defining a \textit{von Neumann winner}: a policy whose expected performance surpasses or matches that of any alternative policy at least 50\% of the time \citep{dudik2015contextual}. A key difference is that preference annotations in the CDB setting are obtained sequentially (online), while learning from human preferences typically relies on a fixed, offline collection of preference-labeled action pairs \citep{yan2022human}. In a similar vein, \textit{preference-based RL} (PbRL) utilizes binary preferences derived from an \textit{unknown} underlying scoring function in place of explicit rewards \citep{BusaFekete2014,ruiz2023dueling}. Existing PbRL algorithms, including those capable of leveraging off-policy preference data, commonly proceed by first estimating the latent reward or scoring function before optimizing against it \citep{jain2013learning,BusaFekete2014,christiano2017deep,sadigh2017active,kupcsik2018learning}. Departing from this two-step paradigm, we propose a unified framework that learns a policy in one stage by directly aligning with observed preferences.

\section{Preliminaries}\label{section:prelims}

We summarize the RLHF process as detailed by \citeauthor{ziegler2020finetuning} (and subsequent works: \citep{stiennon2022learning, bai2022training, ouyang2022training}). This process is generally segmented into three main components: (1) supervised fine-tuning (SFT); (2) the collection of preference data alongside reward model training; and (3) policy improvement via reinforcement learning.

\textbf{SFT}: The typical RLHF workflow starts by fine-tuning a pre-trained language model with supervised learning, using a curated dataset relevant to the downstream task (e.g., dialogue, summarization). This yields an initial model, denoted by $\pi^\text{SFT}$.

\textbf{Reward Modelling Phase}: In the subsequent stage, prompts $x$ are given to the SFT model, which then generates answer pairs $(y_1, y_2)\sim \pi^\text{SFT}(y \mid x)$. These responses are evaluated by human annotators, who select their preferred completion; this preference is indicated as $y_w\succ y_l \mid x$, where $y_w$ and $y_l$ are the chosen and rejected responses, respectively. The preferences reflect judgments informed by an implicit reward function $r^*(y, x)$, which remains unobserved. To model this preference behavior, the Bradley-Terry (BT) model \cite{bradley1952rankanalysis} is frequently employed, though extensions such as the Plackett-Luce ranking models \citep{plackett1975analysis, luce2012individual} are compatible with the methodology when multiple responses can be ranked. According to the BT model, the observed preference distribution $p^*$ is formalized as follows:

\begin{equation}\label{eq:bradley-terry}
    p^*(y_1\succ y_2 \mid x)=\frac{\exp\left(r^*(x, y_1)\right)}{\exp\left(r^*(x, y_1)\right) + \exp\left(r^*(x, y_2)\right)}.
\end{equation}

Given a fixed dataset of human preference comparisons $\mathcal{D}=\bigl\{x^{(i)}, y_w^{(i)}, y_l^{(i)}\bigr\}_{i=1}^N$, assumed to be sampled according to $p^*$, we can introduce a parametric reward function $r_{\phi}(x, y)$ and fit its parameters by maximizing likelihood. When reframing this as a binary classification problem, the associated negative log-likelihood objective becomes:

\begin{equation}\label{eq:reward_model}
    \mathcal{L}_R(r_{\phi}, \mathcal{D}) = -\mathbb{E}_{(x, y_w, y_l)\sim \mathcal{D}}\bigl[\log \sigma(r_{\phi}(x, y_w)- r_{\phi}(x, y_l))\bigr]
\end{equation}

where $\sigma$ denotes the logistic sigmoid. In practice with language models, $r_{\phi}(x, y)$ is commonly initialized from the SFT model $\pi^\text{SFT}(y \mid x)$, supplemented with a linear projection on top of the transformer’s final hidden layer to yield a scalar reward estimate \cite{ziegler2020finetuning}. To regularize the reward estimator and reduce its variance, it is standard practice to normalize outputs so that $\mathbb{E}_{x,y\sim \mathcal{D}}\left[r_\phi(x, y)\right] = 0$ for every $x$.

\textbf{RL Fine-Tuning Phase}: At this stage, the learned reward signal is used to guide the subsequent language model updates. Following prior methodologies~\citep{jaques2017sequence, jaques2020human}, the objective is written as

\begin{equation}\label{eq:RL}
\max_{\pi_{\theta}}  \mathbb{E}_{x\sim \mathcal{D}, y\sim \pi_{\theta}(y \mid x)}\bigl[r_{\phi}(x, y)\bigr] - \beta\mathbb{D}_{\textrm{KL}}\bigl[\pi_{\theta}(y\mid x)\mid \mid \pi_\text{ref}(y\mid x)\bigr],
\end{equation}

where $\beta$ moderates how much the new policy is permitted to diverge from a fixed reference distribution $\pi_\text{ref}$, typically set to the SFT baseline $\pi^\text{SFT}$. In most cases, the policy $\pi_\theta$ is initialized to this reference policy. The KL penalty is crucial for restricting policy drift, ensuring the agent operates in regions of the sample space where the reward estimator remains reliable, while also preserving output variability and avoiding collapse to overly narrow solution sets. Since outputting sequences in language models is a discrete process, direct optimization of this functional is not possible; thus, reinforcement learning strategies are used. A conventional approach \citep{ziegler2020finetuning, stiennon2022learning, bai2022training, ouyang2022training} is to use the modified reward $r(x, y) = r_{\phi}(x, y) -\beta (\log \pi_{\theta}(y\mid x) - \log \pi_\text{ref}(y\mid x))$, which is then maximized using PPO \cite{schulman2017proximal}.

\section{Direct Preference Optimization}\label{sec:DPO}

Considering the substantial difficulties in scaling up RL methods to large models such as LMs, our objective is to construct a straightforward policy optimization method that can utilize preference data without a full RL pipeline. In contrast to previous RLHF frameworks that entail reward learning followed by separate RL optimization, our proposal is to parameterize the reward in such a way that one can analytically determine the corresponding optimal policy, dispensing with the need for iterative RL. As elaborated in the following sections, the central insight is to establish a direct, analytic relationship between the parameterized reward and the optimal policy, allowing us to recast the reward function optimization problem as a direct optimization over policies. Through this reparametrization, we can forego explicit standalone reward modeling, yet remain aligned with human preference models, notably the Bradley-Terry formulation. In this formulation, the

\textbf{Deriving the DPO objective.} Beginning with the standard RL objective as formulated in Eq.~\ref{eq:RL}, and employing a generic reward function $r$, we note that existing research~\citep{peters2007reinforcement, peng2019advantage, korbak2022reinforcement, go2023aligning} demonstrates that the solution maximizing the reward under a KL-divergence constraint is:

\begin{equation}\label{eq:op_policy}
    \pi_r(y\mid x) = \frac{1}{Z(x)}\pi_\text{ref}(y\mid x)\exp\left(\frac{1}{\beta}r(x, y)\right),
\end{equation}

Here, the normalization constant $Z(x) =\sum_{y}\pi_\text{ref}(y\mid x)\exp\left(\frac{1}{\beta}r(x, y)\right)$ acts as the partition function, and detailed steps are provided in Appendix \ref{app:derivation1}. Even when the reward model $r_{\phi}$ is an MLE estimate for the actual reward $r^*$, computation of $Z(x)$ remains costly~\citep{korbak2022reinforcement, go2023aligning}, limiting direct applicability. Nevertheless, Eq.~\ref{eq:op_policy} can be reformulated to present the reward as a function of its optimal policy $\pi_r$, baseline policy $\pi_\text{ref}$, and the unknown partition function $Z(\cdot)$. Taking logarithms on both sides of Eq.~\ref{eq:op_policy} and rearranging yields:

\begin{equation}\label{eq:main_eq}
    r(x,y) =\beta \log \frac{\pi_r(y\mid x)}{\pi_\text{ref}(y\mid x)} + \beta \log Z(x).
\end{equation}

This transformation can be similarly utilized for the ground-truth reward $r^*$ and its optimal policy $\pi^*$. In particular, the Bradley-Terry model is solely dependent on reward differences for competing responses, specifically, ${p^*(y_1 \succ y_2 \mid x) = \sigma(r^*(x, y_1) - r^*(x, y_2))}$. Inserting the new representation from Eq.~\ref{eq:main_eq} for $r^*(x, y)$ into Eq.~\ref{eq:bradley-terry} for the preference model eliminates the partition function, expressing the probability of a human preferring one output over another entirely in terms of the optimal policy $\pi^*$ and the reference policy $\pi_\text{ref}$. Therefore, the Bradley-Terry model implies the optimal RLHF policy $\pi^*$ should satisfy the relationship:

\begin{equation}\label{eq:objective}
    p^*(y_1\succ y_2 \mid x)=\frac{1}{1 + \exp\left(\beta \log \frac{\pi^*(y_2\mid x)}{\pi_\text{ref}(y_2\mid x)} - \beta \log \frac{\pi^*(y_1\mid x)}{\pi_\text{ref}(y_1\mid x)}\right)}
\end{equation}

See Appendix~\ref{app:derivation2} for the full derivation. Although this result leverages the Bradley-Terry model, the same analytical strategy can be employed for the Plackett-Luce family of models~\citep{plackett1975analysis, luce2012individual}, as detailed in Appendix~\ref{app:plackett_luce_models}.

With the human preference likelihood now rewritten in terms of the optimal policy, we are positioned to specify a maximum likelihood estimation objective for a parametric policy $\pi_\theta$. Drawing analogy to approaches based on explicit reward modeling (see Eq.~\ref{eq:reward_model}), this leads us to the following policy objective:

\begin{equation}\label{eq:optimum_model}
    \mathcal{L}_\text{DPO}(\pi_{\theta}; \pi_\text{ref}) = -\mathbb{E}_{(x, y_w, y_l)\sim \mathcal{D}}\left[\log \sigma \left(\beta \log \frac{\pi_{\theta}(y_w\mid x)}{\pi_\text{ref}(y_w\mid x)} - \beta \log \frac{\pi_{\theta}(y_l\mid x)}{\pi_\text{ref}(y_l\mid x)}\right)\right].
\end{equation}

Through this approach, we effectively learn an implicit reward function via a different parameterization, for which the optimal policy aligns with $\pi_\theta$. Additionally, our method corresponds to optimizing a reparametrized Bradley-Terry model, granting it established theoretical guarantees, including consistency provided that the preference data distribution satisfies certain conditions \cite{bong2022generalized}. A deeper examination of DPO's theoretical attributes and how they relate to previous studies appears in Section~\ref{sec:theory}.

\textbf{Mechanistic insight into the DPO update.} To develop an intuition for DPO, it is instructive to inspect the gradient of the loss $\mathcal{L}_\text{DPO}$. This gradient with respect to the parameters $\theta$ is given by:

\begin{multline*}\label{eq:gradient}
    \nabla_\theta \mathcal{L}_\text{DPO}(\pi_\theta;\pi_\text{ref}) = \\ -\beta\mathbb{E}_{(x, y_w, y_l) \sim \mathcal{D}} \left[\underbrace{\sigma(\hat{r}_\theta(x, y_l) - \hat{r}_\theta (x, y_w))}_\text{emphasized when the reward assignment is incorrect}\left[\underbrace{\nabla_\theta\log \pi(y_w \mid x)}_\text{boosting likelihood for $y_w$} - \underbrace{\nabla_\theta\log\pi(y_l \mid x)}_\text{penalizing likelihood for $y_l$}\right]\right],
\end{multline*}

with $\hat{r}_\theta(x, y) = \beta \log \frac{\pi_\theta(y \mid x)}{\pi_\text{ref}(y \mid x)}$ serving as the implicit reward determined by the policy model $\pi_\theta$ in relation to the baseline $\pi_\text{ref}$ (refer also to Section~\ref{sec:theory}). Conceptually, this gradient increases the probability assigned to favorable outputs $y_w$, while diminishing it for unfavorable outputs $y_l$. Crucially, the magnitude of this update depends on how much the implicit reward function $\hat{r}_\theta$ incorrectly rates $y_l$ above $y_w$, modulated by $\beta$ and thus reflecting the misordering strength and the impact of the KL regularization. Empirically, this weighting proves to be significant; in our studies, omitting the weighting factor (i.e., employing a simplified version of the method) leads to detrimental effects such as model degeneration (see Appendix Table~\ref{tab:unlikelihood_generations}).

\textbf{DPO Overview.} The standard workflow for DPO involves the following steps: 1) For each prompt $x$, generate candidate completions $y_1, y_2$ sampled from $\pi_\text{ref}(\cdot \mid x)$ and use human annotations to build a collection of preference pairs, forming the offline dataset $\mathcal{D} = \{x^{(i)}, y_w^{(i)}, y_l^{(i)}\}_{i=1}^N$; 2) Update the parameters of the language model $\pi_\theta$ by minimizing the loss $\mathcal{L}_\text{DPO}$, given the fixed reference policy $\pi_\text{ref}$, the dataset $\mathcal{D}$, and the selected $\beta$. 
In applied settings, preference datasets that have already been constructed—typically from public sources—are preferable to collecting new human judgments and samples. Because such datasets are usually drawn using a model fine-tuned with supervised learning ($\pi^\text{SFT}$), we set $\pi_\text{ref} = \pi^\text{SFT}$ whenever this policy is accessible. If $\pi^\text{SFT}$ is unavailable, $\pi_\text{ref}$ is initialized by fitting it via maximum likelihood estimation on the preferred completions ${(x, y_w)}$, formally: ${\pi_\text{ref} = \argmax_{\pi}\mathbb{E}_{x, y_w \sim \mathcal{D}}\left[\log \pi(y_w \mid x)\right]}$. This approach addresses the issue of distribution mismatch between the ideal, but inaccessible, reference distribution and the practical reference policy used within DPO. Additional information regarding implementation details and hyperparameter selection is provided in Appendix~\ref{app:implementation}.

\section{Theoretical Analysis of DPO} In this section, we delve deeper into the theoretical understanding of the DPO framework, offering interpretations, formal guarantees, and discussing why DPO has certain benefits over actor-critic strategies commonly used in RLHF pipelines (for example, PPO~\cite{schulman2017proximal}).

\label{sec:theory}

\subsection{Your Language Model Is Secretly a Reward Model} DPO eliminates the necessity of explicit reward modeling and reinforcement learning steps by relying on a unified maximum likelihood objective. Importantly, the objective in Eq.~\ref{eq:main_eq} corresponds to the Bradley-Terry model, using a specific reward representation $r^*(x, y) = \beta \log\frac{\pi^*_\theta(y \mid x)}{\pi_\text{ref}(y \mid x)}$. Optimization of the model $\pi_{\theta}$ in this context is mathematically analogous to reward model fitting as in Eq.~\ref{eq:reward_model}, provided an appropriate transformation of variables. The following exposition constructs the theoretical justification for this reparameterization, demonstrates that it does not limit the expressiveness of reward functions that can be learned, and establishes that it facilitates the exact recovery of an optimal policy. To lay the groundwork, we start by formalizing an equivalence relation among reward functions.

\begin{definition}
We define two reward functions $r(x, y)$ and $r'(x, y)$ to be equivalent if and only if there exists a function $f$ such that $r(x, y)-r'(x, y) = f(x)$.
\end{definition}

It is straightforward to verify that this defines an equivalence relation, which partitions the space of reward functions into distinct equivalence classes. The following two lemmas formalize important properties arising from this equivalence:

\begin{lemma}\label{lemma:same_prefrence}
Within the Plackett-Luce framework, including the special case of Bradley-Terry, any two reward functions belonging to the same equivalence class yield identical preference distributions.
\end{lemma}

\begin{lemma}\label{lemma:same_policy}
Any two reward functions in the same equivalence class determine the same optimal policy for the constrained RL setting.
\end{lemma}

The proofs for these statements are elementary and can be found in Appendix~\ref{app:lemma1}. The first lemma reflects a recognized identifiability issue in the Plackett-Luce family \cite{plackett1975analysis}, wherein extra constraints must be imposed to ensure meaningful maximum likelihood estimates as described in Eq.~\ref{eq:reward_model} \cite{bong2022generalized}. The second lemma indicates that within a given equivalence class, all reward functions yield identical optimal policies; therefore, it suffices to identify any representative from the optimal equivalence class when solving for our final objective. In Appendix~\ref{app:thm1}, we provide the proof of the following theorem:

\begin{theorem}\label{thm:main}
Suppose mild regularity conditions hold. Then every reward class compatible with the Plackett-Luce (and in particular, the Bradley-Terry) models can be expressed in the form ${r(x, y) = \beta \log \frac{\pi(y\mid x)}{\pi_\text{ref}(y\mid x)}}$ for some model $\pi(y\mid x)$ relative to a reference distribution $\pi_\text{ref}(y\mid x)$.
\end{theorem}

\begin{sproof}
Let $r(x, y)$ be any reward function and let $\pi_r(y \mid x)$ denote the corresponding optimal model as specified in Eq.~\ref{eq:op_policy}. Our objective is to show that for every reward in the equivalence class of $r$, there is a representation using the aforementioned parameterization. To this end, define the projection $f$ as
\begin{equation}
    f(r; \pi_\text{ref}, \beta)(x, y) = r(x, y) - \beta\log\sum_{y}\pi_\text{ref}(y\mid x)\exp\left(\frac{1}{\beta}r(x, y)\right)
\end{equation}
Here, the operator $f$ normalizes $r(x, y)$ by subtracting the logarithm of the associated partition function, and this term is solely a function of $x$, meaning $f(r; \pi_\text{ref}, \beta)(x, y)$ remains within the equivalence class of $r(x, y)$. Substituting $r$ according to the right-hand side of Eq.~\ref{eq:main_eq}, which is valid for all $r$, yields $f(r; \pi_\text{ref}, \beta)(x, y) = \beta \log \frac{\pi_r(y\mid x)}{\pi_\text{ref}(y\mid x)}$. Therefore, the projection $f$ yields a representative of the equivalence class with the desired functional form, demonstrating that this reparameterization captures all reward functions of interest without loss of generality.
\end{sproof}

Theorem~\ref{thm:main} can also be interpreted as identifying the specific reward function chosen by the DPO reparameterization within each equivalence class, namely, the reward function that fulfills:

\begin{equation}\label{eq:lag_p}
     \sum_{y}\underbrace{\pi_\text{ref}(y\mid x)\exp\left(\frac{1}{\beta}r(x, y)\right)}_{=\pi(y\mid x)\text{, using Thm.~\ref{thm:main} reparam.}} = 1,
\end{equation}

This indicates that $\pi(y\mid x)$ forms a legitimate probability distribution, as all probabilities are non-negative and their total equals one. Examining Eq.~\ref{eq:lag_p} in the context of Eq.~\ref{eq:op_policy} reveals that this equation represents the partition function of the optimal policy defined by the reward $r(x, y)$. The principal contribution of the DPO algorithm is to introduce suitable constraints within the otherwise underdetermined Plackett-Luce family of preference models (with a particular emphasis on the Bradley-Terry model). This allows for maintaining the representational power of the reward model class, while guaranteeing that the optimal policy from Eq. \ref{eq:op_policy} remains explicitly solvable for any input $x$.

\subsection{Instability of Actor-Critic Algorithms} Our proposed framework further enables analysis of the sources of instability observed in commonly used actor-critic methods in RLHF, like PPO. Focusing on the RL fine-tuning phase described in Section \ref{section:prelims}, we draw on connections with the control as inference paradigm \cite{levine2018reinforcement}, as applied to the constrained RL formulation presented in \ref{eq:RL}. We consider a parameterized family of policies $\pi_{\theta}(y\mid x)$, with training objectives to minimize $\mathbb{D}_{\text{KL}}[\pi_{\theta}(y|x) \mid \mid \pi^*(y\mid x)]$, where $\pi^*$ is defined as the optimal policy from Eq.~\ref{eq:optimum_model} for the reward function $r_{\phi}(y, x)$. After algebraic manipulation, this yields the following objective function:

\begin{equation}\label{eq:AC}
    \max_{\pi_{\theta}}\mathbb{E}_{\pi_{\theta}(y\mid x)}\left[\underbrace{r_{\phi}(x, y) -\beta\log\sum_{y}\pi_\text{ref}(y\mid x)\exp\left(\frac{1}{\beta}r_{\phi}(x, y)\right)}_{f(r_{\phi}, \pi_\text{ref}, \beta)} - \underbrace{\beta\log\frac{\pi_{\theta}(y\mid x)}{\pi_\text{ref}(y\mid x)}}_{\text{KL}}\right]
\end{equation}

The objective considered here is identical to that optimized in earlier studies 
\citep{ziegler2020finetuning, stiennon2022learning, bai2022training, ouyang2022training}, where the DPO-equivalent reward is employed for the reward function $r_{\phi}$. Within this framework, the normalization component in $f(r_{\phi}, \pi_\text{ref}, \beta)$ can be interpreted as the soft value function associated with the reference policy $\pi_\text{ref}$. Although this normalization term does not alter the set of optimal policies, omitting it can introduce substantial variance into the policy gradient estimates, potentially destabilizing the training process. One strategy to address the normalization is to learn an explicit value function; however, optimizing such a function can itself pose challenges. Alternatively, previous research has often employed a baseline reflecting human-generated completions, which serves as a single-sample Monte Carlo approximation to the normalization. In contrast, the DPO approach reparameterizes the reward function so that such baselines are unnecessary.

\section{Experiments}
This section presents an empirical assessment of DPO's capacity to learn policies directly from preference data. We begin with a controlled experiment in text generation, investigating how effectively DPO balances reward maximization against minimizing KL-divergence from a reference policy when compared to established algorithms for preference optimization like PPO. Subsequently, we extend our analysis to larger-scale models and more complex RLHF settings, encompassing tasks such as summarization and dialogue. Our results indicate that, with minimal hyperparameter tuning, DPO frequently matches or surpasses the performance of strong baselines including RLHF with PPO and the selection of the highest-reward trajectory among $N$ samples. The setup of our experiments is detailed prior to the main findings, with further specifics available in Appendix~\ref{app:exp_details}.

\textbf{Tasks.} We investigate three distinct tasks within the domain of open-ended text generation. Across all tasks, each algorithm derives its policy from a preference dataset $\mathcal{D}=\bigl\{x^{(i)}, y_w^{(i)}, y_l^{(i)}\bigr\}_{i=1}^N$. In the \textbf{controlled sentiment generation} setting, $x$ refers to an initial segment of a movie review from the IMDb dataset \cite{maas-EtAl:2011:ACL-HLT2011}, with the objective of generating a continuation $y$ that exhibits positive sentiment. To facilitate an objective evaluation, we create preference pairs for generated continuations using a sentiment classifier trained in advance, where $p(\text{positive}\mid x,y_w)>p(\text{positive}\mid x,y_l)$. For SFT, GPT-2-large is further trained on the IMDb training split until the model converges, as elaborated in App~\ref{app:sentiment_details}. The \textbf{summarization} task uses $x$ as a Reddit forum post, requiring the policy to produce a concise summary $y$ capturing the post's main ideas. We adopt the Reddit TL;DR summarization dataset \citep{volske-etal-2017-tl} and integrate human preference data obtained by \citeauthor{stiennon2022learning}. For SFT, we utilize a model fine-tuned on human-composed summaries of forum posts\footnote{\url{https://huggingface.co/CarperAI/openai_summarize_tldr_sft}}, and leverage the TRLX \citep{leandro_von_werra_2023_7790115} library to implement RLHF. The dataset with human preferences, created by \citeauthor{stiennon2022learning}, is based on samples drawn from a comparably trained SFT model. For the \textbf{single-turn dialogue} task, $x$ denotes a user-issued prompt, which may encompass inquiries from astrophysics to personal advice. The goal for the policy is to generate an informative and engaging reply $y$. We employ the Anthropic Helpful and Harmless dialogue dataset \citep{bai2022training}, containing 170,000 conversational transcripts between a person and an automated assistant; each transcript concludes with two responses generated by a large (unspecified) language model and a preference indicator highlighting the preferred response. In this scenario, as no specialized SFT model exists, we train an off-the-shelf language model exclusively on preferred responses, creating the SFT baseline.

\textbf{Evaluation.} We employ two complementary strategies for evaluation in our experiments. To assess the capacity of each algorithm to maximize the constrained reward objective, we operate in a controlled sentiment generation framework where the evaluation measures the trade-off between reward attained and KL-divergence from the reference policy. Here, the full trade-off curve can be characterized since access to the true reward function—a sentiment classifier—is available. By contrast, outside the controlled setting, the exact reward function is unknown. Thus, for summarization and single-turn dialogue, we evaluate each algorithm via its \textit{win rate} against a baseline policy, using GPT-4’s outputs as stand-ins for human judgments of summary quality and dialogue helpfulness. In the case of summarization, the baseline consists of test set reference summaries, while for dialogue, the baseline is the preferred response as indicated in the test data. Prior work indicates large language models may serve as better automated evaluators than commonly used metrics \citep{Chen2023ExploringTU}; however, in Section~\ref{sec:human-judgments}, we further substantiate our use of GPT-4 through a dedicated human study. Results show that GPT-4’s decisions exhibit a strong correlation with human ratings, with agreement levels at least as high as, or higher than, those found between pairs of human annotators.

\textbf{Methods.} Alongside DPO, we benchmark several prevalent methods for aligning language models with human preferences. For summarization, we examine zero-shot prompting using \textbf{GPT-J} \citep{gpt-j}, whereas in dialogue we deploy 2-shot prompting with \textbf{Pythia-2.8B} \citep{biderman2023pythia}. Further, we assess a standard \textbf{SFT} model and introduce \textbf{Preferred-FT}, which undergoes supervised fine-tuning on the preferred completions $y_w$: these are generated either from the SFT model (for controlled sentiment and summarization tasks) or a generic language model (for dialogue). The \textbf{Unlikelihood} approach~\citep{welleck2019neural} constitutes another pseudo-supervised baseline, wherein the model is optimized to increase the likelihood of $y_w$ while simultaneously reducing the probability of $y_l$; the `unlikelihood' term can be weighted by a coefficient $\alpha\in[0,1]$. Reinforcement learning is considered through \textbf{PPO} \citep{schulman2017proximal}, using rewards learned from preference data, and \textbf{PPO-GT}, an oracle variant trained directly on the ground-truth reward available in controlled sentiment generation. We employ two PPO-GT implementations in sentiment experiments: a standard public version \cite{leandro_von_werra_2023_7790115} and a custom variant with reward normalization and tailored hyperparameters for enhanced results; these adjustments are likewise applied in PPO runs with learned rewards. As an additional strong—but computationally demanding—baseline, we utilize \textbf{Best of $N$}: the system samples $N$ completions per input from the SFT model (or Preferred-FT for dialogue) and selects the response that obtains the highest score as determined by a learned reward model. Although this approach can reach high performance by separating reward model quality from PPO optimization, its infeasibility at scale stems from the requirement to generate $N$ completions for each test prompt.

\subsection{How well can DPO optimize the RLHF objective?}

In RLHF algorithms, the reward maximization objective typically incorporates a KL penalty, constraining the policy from diverging excessively from the reference policy while promoting reward acquisition. Accordingly, algorithm comparisons should account for both attained reward and the associated KL divergence, since surpassing the reward at the expense of greatly increased KL divergence is suboptimal. In Figure~\ref{fig:frontier-tldr-main}, we plot the relationship between reward and KL divergence across different algorithms on the sentiment benchmark. For each algorithm, we conduct several training sessions, varying the policy conservativeness parameter each time (using target KL values $\in\{3,6,9,12\}$ for PPO, $\beta \in \{0.05,0.1,1,5\}$, $\alpha\in\{0.05,0.1,0.5,1\}$ for unlikelihood, and random seeds for preferred-FT), yielding 22 experiments in total. Following every 100 training updates until convergence, each model is evaluated on a designated test set. We record both the mean reward according to the true reward model and the mean sequence-level KL\footnote{This refers to the aggregate of KL-divergence computed at each timestep.} with respect to the reference policy, $\text{KL}\left(\pi\mid \mid \pi_\text{ref}\right)$. Our experiments indicate that DPO yields the most favorable reward-KL tradeoff, offering the highest rewards for relatively low KL values. This outcome is particularly striking for several reasons: first, while DPO and PPO are tasked with optimizing the same objective, DPO demonstrates superior sample efficiency, achieving a reward/KL Pareto frontier that surpasses PPO's. Second, DPO outperforms PPO not only in general, but also in settings where PPO is granted access to ground-truth rewards (PPO-GT).

\subsection{Can DPO scale to real preference datasets?}
\label{sec:dpo-real-datasets}
We proceed to assess the effectiveness of DPO when fine-tuned on tasks such as summarization and single-turn dialogue. It is well-established that for summarization, standard automated metrics like ROUGE exhibit weak alignment with actual human preferences~\citep{stiennon2022learning}. Previous research indicates that training language models with PPO guided by human preferences enhances summary quality. To benchmark various techniques, we generate outputs from the test partition of the TL;DR summarization dataset and measure the mean win rate relative to reference completions. Sampling is performed across a range of temperatures (0.0 to 1.0) for all models, with win rates presented in Figure~\ref{fig:frontier-tldr-main} (right). All methods, including DPO, PPO, and Preferred-FT, further train a common GPT-J SFT model\footnote{\url{https://huggingface.co/CarperAI/openai_summarize_tldr_sft}}. Results show that DPO attains an approximate win rate of 61\% at temperature 0.0, outperforming PPO, which reaches about 57\% at its own best temperature setting. Additionally, DPO surpasses the highest win rate achieved by the best of $N$ baseline. It is important to mention that no substantial tuning was performed on DPO's $\beta$ hyperparameter, implying that these findings may not fully capture DPO’s capabilities. Furthermore, DPO demonstrates greater resilience to variations in sampling temperature, whereas PPO’s win rate deteriorates at higher temperatures, approaching that of the underlying GPT-J model. Preferred-FT yields little to no improvement relative to the SFT baseline. In further analysis presented in Section~\ref{sec:human-judgments}, human raters preferred samples from DPO at temperature 0.25 over those from PPO at the same temperature 58\% of the time.

For single-turn dialogue evaluation, we examine the various approaches on the subset of the Anthropic HH dataset’s test split \citep{bai2022training} that contains just one exchange between human and assistant. In the GPT-4 assessment, we use the completions labeled as preferred within the test set as references, and measure each method’s win rate accordingly. Given the absence of a standard SFT model for this scenario, our process begins with a pre-trained Pythia-2.8B model; we employ Preferred-FT to train a reference model on the selected preferred completions, ensuring these outputs remain in-distribution, before subsequently training with DPO. As part of our comparison, we include the best result among 128 Preferred-FT completions (our observations indicate the Best of $N$ approach saturates at 128 completions for this benchmark; refer to Appendix Figure~\ref{fig:best-of-n}), alongside a 2-shot prompted Pythia-2.8B base model. In all cases, DPO achieves equal or superior results relative to these baselines at optimal sampling temperatures for each method. Additionally, we assess an RLHF model—trained with PPO on the Anthropic HH dataset \footnote{\url{https://huggingface.co/reciprocate/ppo_hh_pythia-6B}} and provided by a recognized source \footnote{\url{https://github.com/CarperAI/trlx/tree/main/examples/hh}}—but we are unable to identify any prompting or temperature configuration that outperforms the base Pythia-2.8B model. Drawing on our TL;DR findings and the shared reward function optimized by both methods, we take Best of 128 to be a practical benchmark representing PPO-level performance. Summarizing our findings, DPO emerges as the only method that is both computationally efficient and consistently surpasses the preferred completions in the Anthropic HH dataset, matching or exceeding the computationally intensive Best of 128 baseline. Lastly, Figure~\ref{fig:dialogue-main} illustrates that DPO attains peak performance in relatively few iterations.

\subsection{Generalization to a new input distribution}

To more thoroughly assess how PPO and DPO perform when exposed to out-of-distribution data, we test the policies derived from our Reddit TL;DR summarization task on a separate corpus: the test set of the CNN/DailyMail dataset, comprising news articles \citep{nallapati-etal-2016-abstractive}. We utilize the optimal sampling temperatures previously identified for TL;DR (0 and 0.25) in this evaluation. Table~\ref{tab:ood} displays the outcomes. For assessment, we calculated the win rates of GPT-4 versus the reference summaries in each dataset, employing the same GPT-4 (C) prompt as in the Reddit TL;DR study, with the sole adjustment of substituting ``forum post'' with ``news article''. On this novel dataset, DPO maintains a substantial performance lead over PPO. These findings offer preliminary support that DPO policies are capable of generalizing to new distributions as effectively as those trained with PPO, despite not utilizing the extra unlabeled Reddit TL;DR prompts leveraged by PPO.

\subsection{Validating GPT-4 judgments with human judgments}
\label{sec:human-judgments}

To assess the fidelity of GPT-4’s ratings, we conducted a human evaluation centered on the TL;DR summarization task, incorporating two distinct GPT-4 prompting strategies. The \textbf{GPT-4 (S)} (simple) prompt requests a choice for the summary that better captures the key content from the post. The \textbf{GPT-4 (C)} (concise) prompt additionally solicits a judgment regarding conciseness; this prompt is considered because under the \textbf{GPT-4 (S)} prompt, GPT-4 demonstrates a bias towards longer, more repetitive summaries than those favored by human evaluators. Appendix~\ref{app:prompts} provides full prompt details. Three pairwise comparisons are carried out: the best-performing method (DPO, temp. 0.25), the worst-performing (PPO, temp. 1.0), and an intermediate baseline (SFT, temp. 0.25), each contrasted against greedily-sampled PPO (the optimal temperature for PPO), to span a broad range of summary quality. Across both prompting strategies, agreement between GPT-4 and humans is approximately on par with inter-annotator agreement among human raters themselves, indicating that GPT-4 can act as an effective surrogate for human preference assessments (for resource reasons, multiple human ratings are only obtained for the DPO and PPO-1 comparisons). Notably, the \textbf{GPT-4 (C)} prompt yields win rates most closely aligned with human evaluations; thus, we rely on this prompt for primary results in Section~\ref{sec:dpo-real-datasets}. Comprehensive details on the human evaluation—including rater interface and volunteer information—are available in Appendix~\ref{app:human-study}.

\section{Discussion}
Preference-based learning offers a robust and flexible approach for developing language models that are both competent and aligned with human values. In this work, we present DPO, a straightforward training strategy for learning from preferences that circumvents the complexities of reinforcement learning. DPO bypasses the necessity of adapting the preference learning task to fit conventional RL methodologies, opting instead to establish a direct relationship between language model policies and reward functions. This alignment allows for training models to reflect human preferences through a standard cross-entropy objective, dispensing with reinforcement learning altogether and without sacrificing generality. Remarkably, DPO matches or outperforms state-of-the-art RLHF approaches, such as those based on PPO, all while requiring minimal hyperparameter adjustment. As a result, DPO significantly lowers the hurdles to training language models using human preference data.

\textbf{Limitations \& Future Directions.} This study opens several avenues for subsequent investigation. One area concerns how DPO-trained policies extrapolate to out-of-distribution settings compared to those trained via explicit reward optimization. Preliminary experiments imply that the generalization capabilities of DPO are on par with PPO-based methods, yet more extensive evaluation is necessary. For instance, it remains an open question whether utilizing self-labeling via the DPO policy is comparably effective for leveraging prompts that lack labels. Another important direction is understanding the manifestation of reward over-optimization in the context of direct preference optimization, and whether the minor performance drop observed in Figure~\ref{fig:dialogue-main}-right can be attributed to this phenomenon. Additionally, while current experiments encompass models with up to 6B parameters, exploring the feasibility and performance of DPO when applied to models of significantly larger scale represents a compelling path forward. In terms of model assessment, our findings suggest that GPT-4's computed win rates are influenced by prompt formulation; refining methods for extracting reliable evaluations from automated systems remains a key research question. Finally, DPO's methodology has the potential to extend beyond language modeling to a wide array of generative modeling tasks across different modalities.